\documentclass[11pt]{article}
\usepackage[T1]{fontenc}
\usepackage{lmodern}
\usepackage[letterpaper,margin=0.86in,top=0.72in,bottom=0.77in]{geometry}
\usepackage{amsmath,amssymb,xcolor,hyperref}
\definecolor{ink}{HTML}{183E55}
\hypersetup{colorlinks=true,urlcolor=ink,pdftitle={Sharper bounds for adaptive threshold estimation}}
\setlength{\parindent}{0pt}
\setlength{\parskip}{5pt plus 1pt minus 1pt}
\setlength{\abovedisplayskip}{6pt}
\setlength{\belowdisplayskip}{6pt}
\raggedbottom
\newcommand{\E}{\mathbb E}
\newcommand{\Aff}{\operatorname{Aff}}
\newcommand{\hd}{\mathfrak h}
\newcommand{\heading}[1]{{\LARGE\bfseries\color{ink} #1}\par\smallskip}
\newcommand{\src}[2]{\href{https://github.com/davidiach/adaptive-threshold-estimation/blob/97a501383083cc0fa64cb96405444048b8fd4053/#1}{#2}}
\begin{document}
{\small\color{ink} ADAPTIVE THRESHOLD ESTIMATION \hfill MATHEMATICAL REVIEW \quad 17 NOVEMBER 2026}\par
\heading{A stronger constant for the main RMS law}
For every positive-gap law $\mu$ on $(0,1/2]$ and expected flip budget $B\ge1$, let
$h=B^{-1/2}$, $f(d)=\min(1,h/d)$ and $\omega_\mu(h)=\int f\,d\mu$.
In the ordinary hidden-magnitude experiment the minimax RMS risk satisfies
\[
   R_\mu(B)\ \ge\ \frac{4\sqrt{21}-15}{25}\,\omega_\mu(h)
          \ >\ 0.1332\,\omega_\mu(h).
\]
The published coefficient is $1/64$. A four-flip example below puts an upper bound of $1/6$ on the best possible universal coefficient. The results in this note are paper derivations based on the linked statements at commit \texttt{97a5013}.

\textbf{Proof.} Use the two sign laws $q_\pm(d)=1/2\pm f(d)/5$. At the midpoint, conditional on the gap $d$, the hidden sign $S\in\{-1,1\}$ is fair. For a one-coin policy, let $M_n$ be the posterior mean of $S$ given $d$, the independent seed and the first $n$ flips of its potential stream. For $z=M_n$, after observing a flip $Y\in\{-1,1\}$ the new mean is $(z+2dY)/(1+2dzY)$ and $\Pr(Y=y)=(1+2dzy)/2$. Consequently
\[
 \E[(M_{n+1}-M_n)^2\mid\mathcal F_n]
  =\frac{4d^2(1-z^2)^2}{1-4d^2z^2}\ \le\ 4d^2.
\]
When the sign is already known the increment is zero, including at $d=1/2$.
For each bounded local stop $T$, summing the conditional variances from $M_0=0$ gives $\E[f(d)^2 M_T^2]\le4h^2\E T$.

Write $\hd(P,Q)=1-\Aff(P,Q)$. Under the midpoint, the endpoint likelihoods on the stopped transcript augmented by $d$ are $1\pm u$, where $u=2 f(d) M_T/5$. Erasing $d$ can only decrease $\hd$. Since $|u|\le2/5$, the ordinary local transcript laws satisfy
\[
 \hd(P_-,P_+)\le \E[1-\sqrt{1-u^2}]
 \le \frac{\E[u^2]}{1+\sqrt{21}/5}
 \le 4(1-\sqrt{21}/5)\,h^2 \E T.
\]
The manuscript's midpoint work decomposition sums this estimate over coins even with adaptive revisits. The midpoint is admissible, so its expected work is at most $B$. Apply the decomposition at every finite total cap and pass to the stopped transcript by almost-sure termination. Its affinity is at least $4\sqrt{21}/5-3$. The two targets differ by $2\omega_\mu(h)/5$. The squared-loss affinity inequality $R\ge|\Delta\theta|\,\Aff/2$ gives the result.

\textbf{Ceiling on the constant.} Take $\mu=\delta_{1/2}$ and $B=4$. One flip reveals each fresh sign, so four fresh coins give $X\sim{\rm Bin}(4,q)$. The estimator $(X+1)/6$ has MSE $[4q(1-q)+(1-2q)^2]/36=1/36$ for every $q$, and $\omega_\mu(1/2)=1$.

{\small Sources: \src{paper/v6/core.tex}{main theorem and experiment},
\src{paper/v6/appendix-transfer-foundations.tex}{midpoint work decomposition},
\src{paper/v6/appendix-dual.tex}{squared-loss testing by affinity}.}
\newpage
\heading{A smaller cap for the stopping certificate}
Let $P_N$ be the ordinary stopping-program value with at most $N$ flips per coin, $G_N$ the optimized dual with a terminal reveal of the latent state, and $P_\infty$ their common limit. For every positive-gap law, $0<\rho\le1/2$, $B>0$, and integer $N\ge1$, set $t=\sqrt{B/N}$. Then
\[
 G_N\ \le\ P_\infty\ \le\ P_N\ \le\ kG_N,\qquad
 k=\frac{\sqrt{4+t^2}+t}{2}.
\]
Thus $k=1+t/2+O(t^2)$ as $t\to0$. The published joint bound has cube-root excess. The factor-two test $P_N\le2G_N$ now holds for $N\ge\lceil4B/9\rceil$, compared with the published sufficient cap $\lceil4B/3\rceil$.

\textbf{Proof.} Fix admissible endpoint, variance and work priors with finite ordinary information/work ratio $K$ at cap $N$. Mixing both reference priors toward the endpoint midpoint with weight $0<e<1$ preserves the target gap $\Delta$. The manuscript's regularization estimate bounds the resulting genie price by
\[
 K^+(e)\le\max\{K/(1-e)^2,\,4/(Ne^2)\}.
\]
Put $x=\sqrt{BK}/2$. For $x>0$ choose $e=t/(x+t)$, which equalizes the two terms and gives $BK^+(e)/4\le (x+t)^2$. Since $k-1/k=t$, the required denominator comparison follows from the exact identity
\[
 k^2(1+x^2)-1-(x+t)^2=(k^2-1)(x-1/k)^2\ge0.
\]
Hence $\Delta_+^2/(4+BK)\le k^2 G_N^2$. At $K=0$ take $e\uparrow1$ in the bound for this same optimized value. Infinite $K$ gives a zero ordinary quotient. Taking the supremum over ordinary quartets proves the result. The condition $k\le2$ is exactly $t\le3/2$.

\textbf{A lower obstruction for the optimized genie.} For each fixed $B,N$, take $\rho=1/2$ and a single gap $d>0$. In the infinite ordinary dual choose $q_L=0$, $q_H=q_v=q_w=1/2$. Its stopped signed likelihood score is the posterior sign mean $M_T$ on page 1. The variance calculation bounds its information/work ratio by $4d^2$ at every cap, with equality after one flip. Hence $P_\infty\ge (1/2)/\sqrt{4+4Bd^2}$.

For \emph{every} genie quartet, the policy that always buys $N$ flips and receives the sign has information/work at least $4\Delta^2/N$, since $q_v(1-q_v)\le1/4$. This includes singular references under the closed-perspective convention. As $|\Delta|\le1/2$, the optimized value satisfies $G_N\le(1/2)/\sqrt{4+B/N}$. Letting $d\downarrow0$ through admissible positive gaps proves
\[
 \sup_{\mu,\,0<\rho\le1/2} P_\infty/G_N \ \ge\ \sqrt{1+B/(4N)}.
\]
The worst-case relative excess therefore lies between a lower bound of order $B/N$ and an upper bound of order $\sqrt{B/N}$ as $B/N\to0$. The sharp exponent remains open.

{\small Sources: \src{paper/v6/appendix-cap-rate.tex}{published cap theorem},
\src{lean/ThresholdHC/GenieRate/Flow.lean}{formal mixed-price inequality},
\src{paper/v6/appendix-genie.tex}{optimized dual definitions and brackets}.}
\newpage
\heading{A tighter certificate for minimax risk}
In the same rare-event model let $R$ be the minimax RMS risk with expected budget $B$. The infinite stopping program satisfies
\[
       R\ \le\ P_\infty\ \le\ \frac{3\sqrt7}{2}\,R
       \qquad (3\sqrt7/2 \simeq 3.969).
\]
The published upper factor is $8$. Combining this bound with page 2 gives $R\le P_N\le 3\sqrt7\,R<8R$ whenever $N\ge\lceil4B/9\rceil$. A verified feasible upper bound $U$ and genie lower bound $L$ with $U\le2L$ give the same factor for $U$.

\textbf{Proof.} Take an admissible quartet $(q_L,q_H,q_v,q_w)$ in the infinite-cap dual, with nonnegative target gap $\Delta$ and finite ordinary price $K$. For $0<a<1$, let $b=(1-a)/2$ and replace the endpoints by
\[
 q'_- = a q_L+bq_v+bq_w,\qquad
 q'_+ = a q_H+bq_v+bq_w.
\]
They are admissible by convexity and have target separation $a\Delta$. For any capped local policy its transcript laws are affine in the one-coin prior. With densities relative to a common dominating measure, both new densities are at least $b p_v$, their difference is $a(p_H-p_L)$, and the new midpoint work $W_m$ is at least $b W_w$. Consequently
\begin{align*}
 \hd(P'_-,P'_+)
 &=\tfrac12\int \frac{(p'_+-p'_-)^2}
                         {(\sqrt{p'_+}+\sqrt{p'_-})^2}\\
 &\le \frac{a^2}{8b}\int\frac{(p_H-p_L)^2}{p_v}
 \ \le\ \frac{a^2 K}{8b}W_w
 \ \le\ \frac{a^2 K}{8b^2}W_m .
\end{align*}
The integral uses the closed perspective at zero. Finite $K$ forces zero contrast on a null reference cell. A local policy that sometimes takes zero flips follows by conditioning on activation, which depends only on its independent seed.

Apply the midpoint work decomposition to bounded population prefixes, then pass to the limit as on page 1. Writing $x=BK$, the population affinity is at least $1-a^2x/[2(1-a)^2]$. When $x>0$ take $a=(1+\sqrt{3x/2})^{-1}$. This gives affinity at least $2/3$, hence
\[
 R\ \ge\ \frac{\Delta}{3(1+\sqrt{3x/2})}
 \ \ge\ \frac{2}{3\sqrt7}\,\frac{\Delta}{\sqrt{4+x}}.
\]
The last step is Cauchy--Schwarz on $(2,\sqrt{x})$ and $(1/2,\sqrt{3/2})$. At $x=0$ let $a\uparrow1$ to obtain $R\ge\Delta/2$. An infinite price contributes zero to the dual. Taking the supremum over quartets gives the upper comparison, and the published estimator realization gives $R\le P_\infty$.

{\small Sources: \src{paper/v6/appendix-dual.tex}{reference absorption and affinity testing},
\src{paper/v6/appendix-genie.tex}{infinite-cap dual and published factor eight},
\src{paper/v6/appendix-transfer-foundations.tex}{local-to-population transfer}.}
\end{document}
